\subsubsection{Adequacy, grouped holdout, and advisory relaxation}
\label{sec:reduced-od-validation}

Phase~12 separates three questions that must not be conflated.  Full-data
\emph{adequacy} asks how well a fitted model reproduces the observations used
to estimate it.  Grouped \emph{holdout validation} asks how well a model refitted
without selected observations predicts those untouched observations.  An
\emph{advisory relaxation score} identifies a possible next model but neither
constructs nor fits that model automatically.  The implementation is divided
between \path{validation.py}, \path{diagnostics.py}, and \path{lineage.py} in
the reduced-OD inference package.

The minimal objective now accepts an optional immutable calibration mask.  Its
default remains all measurements, so Phases~9--11 are unchanged.  During a
holdout fit, likelihood contributions outside the calibration mask are replaced
by zero before summation.  Predicted means are nevertheless returned for every
measurement row.  Consequently changing holdout counts cannot change the
calibration objective, gradient, or estimate; a regression test changes held-out
counts by 10,000 and obtains exactly the same fitted raw parameters.

Holdout construction is deterministic from a measurement identity, seed, and
group labels.  Supported grouping units include complete vehicle journeys,
stop time series, lines, directions, time blocks, and explicit groups.  Every
group lies wholly in calibration or holdout, both subsets must be nonempty, and
their masks are complementary.  In particular, boarding and alighting rows
from one held-out vehicle journey cannot leak individually into calibration.
Calibration and holdout summaries report their own counts, observed and modeled
totals, likelihood, Poisson and negative-binomial deviances, MAE, RMSE, and
variance-weighted RMSE.  For a Poisson model the negative-binomial deviance is
reported at its Poisson limit; a fitted negative-binomial model uses its fitted
dispersion.

Full-data adequacy reports raw and standardized residuals, threshold exceedance
rates, and grouped summaries for available measurement type, line, direction,
stop, period, vehicle journey, origin zone, destination zone, and transfer
place labels.  These are calibration diagnostics and the report refuses a
masked problem rather than presenting calibration fit as predictive evidence.
The Hessian of the small raw-parameter objective supplies minimum and maximum
curvature eigenvalues and a stabilized condition number.  Near-singular
curvature, a large condition number, and too few observations per fitted
parameter generate explicit weak-identification warnings.

Residual patterns may recommend exactly one of four conventional children:
J1, broad-period cost effects; J2, coarse destination-zone attraction; J3,
coarse origin-zone production; or J4, selected transfer-place correction.  A
candidate is applicable only when its declared grouping metadata contains at
least two groups.  Its score is based on the largest grouped mean standardized
residual and is reduced when observation support per added parameter is weak.
The recommendation object is permanently labeled \texttt{advisory\_only}; it
does not alter J0, add parameters to the Phase-9 objective, select a model, or
claim predictive improvement.  Implementing and estimating these children is
a later model-extension decision.

Lineage nevertheless prepares that decision safely.  A fitted J0 root has a
fingerprinted identity, parameter names, and estimates.  Each proposed J1--J4
child records its parent, appends only explicitly named zero-effect parameters,
and is accepted as a warm start only after the caller demonstrates that the
child prediction at those zeros reproduces the parent prediction within a
declared tolerance.  Such nodes remain marked unfitted.  Thus every proposed
child is auditable without silently turning a diagnostic recommendation into an
estimation run.

On the public 20,000-cell, 2,000-measurement benchmark, cold adequacy including
the first Hessian compilation took about 2.18 seconds and the warm repeated
adequacy calculation about 0.046 seconds.  Four advisory candidates were scored
in about 0.00018 seconds, and a vehicle-journey-grouped 20-percent holdout refit
took about 0.153 seconds after the relevant JAX shapes had been compiled.  These
diagnostic costs remain separate from both preprocessing and detailed-assignment
validation.

